\chapter{The survey: results}\label{chap:04-sondage-resultats}

This chapter gives all the results of the survey: the samples obtained
(Section~\ref{sec:sr-echantillons}), the main comparison with the same grid on both sides, by
criterion (Section~\ref{sec:sr-principale}) then by flaw (Section~\ref{sec:sr-failles}), what
withstands the robustness analyses (Section~\ref{sec:sr-robustesse}), our reading of 2018 against
Pineau's figures (Section~\ref{sec:sr-pineau}), the first 2026 draw and the papers on language models
(Section~\ref{sec:sr-premier}), the trajectory from 2019 to 2024 (Section~\ref{sec:sr-trajectoire}),
the agreement between readers (Section~\ref{sec:sr-accord}), and finally analyses A5 and A6
(Sections~\ref{sec:sr-declarations} and~\ref{sec:sr-filtre}). Every number is taken from the JSON
files of \file{study/} named in each caption.

\section{The samples obtained}\label{sec:sr-echantillons}

\paragraph{The first 2026 draw}
The 50 papers drawn from the 1,252 RL papers of 2025 and 2026 break down as follows: 27 apply what
they call RL to large language models (\file{llm} family), 14 belong to classical RL (\file{classic})
and 9 are theoretical (\file{theory}), 5 of which have no experiment at all. The percentages are
computed over the 45 papers that have experiments.

\paragraph{Our reading of 2018}
The 50 papers from 2018 (21 from ICML, 20 from NeurIPS, 9 from ICLR) include 48 papers with
experiments. Of these 48 papers, one (C14) compares an exact method with no hyperparameter: Q3 and Q6
are \file{na} for it, hence a denominator of 47 for these two criteria. Q4 and Y5 are \file{na} when a
paper compares no baseline.

\paragraph{Classical RL in 2026}
The 50 papers bring together the 14 \file{classic} papers of the first draw and the first 36
\file{classic} papers of the A3 queue (\file{extension\_stats.json}, field \file{samples}). All 50 have
experiments. Seventy papers of the queue were read: reaching the thirty-sixth \file{classic} paper took
57 readings, 18 \file{llm} papers and 3 \file{theory} papers among them despite the title screen.
Thirteen more papers were read after it.

\paragraph{\enquote{50 against 48}}
The blog article and our figures speak of \enquote{50 classical RL papers from 2026} against
\enquote{50 papers from 2018}. Both samples do hold 50 papers read. The 2018 percentages are
nevertheless computed over 48 papers, sometimes 47, because two 2018 papers have no experiment. Those of
2026 are computed over 50 papers, sometimes 49 when a criterion does not apply. The exact denominator of
each share appears in the tables of this chapter. The label \enquote{50 against 50} thus describes the
samples drawn, the label \enquote{48 against 50} the samples analysed for most criteria.

\section{The main comparison, with the same grid}\label{sec:sr-principale}

\paragraph{The two samples and their readers}
The main comparison sets our 50 classical RL papers from 2026 against our reading of 50 papers from
2018, with the same grids and the same harmonised rules (Section~\ref{sec:sm-verdict}). The 2018 papers
and the 2026 papers were read by different readers (\file{study/survey\_extra.json}, field
\file{readers\_year}). An effect peculiar to one year's readers therefore cannot be told apart from the
year itself in this comparison (Section~\ref{sec:sm-annees}). Deviation 36 separates the two by having
the 98 papers reread by eight readers who each read both years (Section~\ref{sec:sr-annees}).
Criterion X8, further on, already escaped this confound.

Table~\ref{tab:sondage-pineau} gives the criteria of grid v1, Table~\ref{tab:sondage-extension} those
of grid v2. For each criterion, the tables give the share in 2018 and in 2026 with their Wilson
intervals, the difference with its Newcombe interval, the $p$-value of Fisher's exact test and the value
adjusted by the Holm correction, applied separately to the 11 v1 criteria and to the 12 v2 criteria.
Figure~\ref{fig:forest-criteres} draws these differences criterion by criterion. Figure~\ref{fig:grille-egale} takes up the eight practices of the blog article.

\gentable{sondage-pineau}

\gentable{sondage-extension}

\fig{forest-criteres}{All the reading criteria, grid v1 (Q, X) and grid v2 (Y), in our reading of 2018
(48 papers with experiments) and in our 50 classical RL papers from 2026, read with the same grids and the
same rules. For each criterion, the two dots give the two shares, the bar the 95\% Newcombe
interval of their difference. The criteria whose difference passes the Holm correction are highlighted. For X3, a
lower share is progress. Data: \file{study/extension\_stats.json}.}

\fig{grille-egale}{Difference, in percentage points, between 50 classical RL papers from 2026 and 50 papers
from 2018 (48 with experiments), read with the same grids and the same rules, for the eight practices
of the blog article. The dot gives the observed difference, the bar its 95\% Newcombe interval. The four
practices in green pass the Holm correction: linking the code (Q5), showing the effect of a setting
(Y9), quantifying the compute (X7), reporting an uncertainty (X2). The four others, the statistical
test (Q8) among them, have intervals that contain zero. Data: \file{study/extension\_stats.json}.}

\paragraph{Four practices changed}
Four differences pass the Holm correction.
\begin{itemize}
  \item \textbf{Code linked (Q5)} goes from 31\% (15 of 48) to 74\% (37 of 50): $+43$ points,
  $[+23,\,+58]$, $p_{\text{Holm}} = 0.0004$.
  \item \textbf{Sensitivity to a hyperparameter shown (Y9)} goes from 44\% (21 of 48) to 80\% (40 of
  50): $+36$ points, $[+17,\,+52]$, $p_{\text{Holm}} = 0.004$.
  \item \textbf{Compute quantified (X7)} goes from 27\% (13 of 48) to 58\% (29 of 50): $+31$ points,
  $[+11,\,+47]$, $p_{\text{Holm}} = 0.02$.
  \item \textbf{Uncertainty reported (X2)} goes from 63\% (30 of 48) to 88\% (44 of 50): $+26$ points,
  $[+8,\,+41]$, $p_{\text{Holm}} = 0.04$. This verdict is the most fragile
  (Section~\ref{sec:sr-robustesse}). X2 counts any error bar, whatever variability it measures: criterion
  X8, further on, asks which.
\end{itemize}
A fifth criterion also passes the correction: the use of neural networks (Q2), from 79\% (38 of 48) to
98\% (49 of 50), $+19$ points, $[+6,\,+32]$, $p_{\text{Holm}} = 0.03$. Q2 describes the population:
classical RL in 2026 is almost entirely deep RL, which it was not quite in 2018. The criterion measures
no reproducibility flaw.

\paragraph{An error bar over what? Criterion X8}
X2 is \file{yes} as soon as a paper reports an uncertainty. Experiment B1 (\cref{sec:b1}) shows, however, that
a bar computed over the episodes of one training run is the trap that produces the 2018 $t$ of 9: it ignores
the variability between random seeds, the one that decides whether two algorithms differ (flaw F11). A paper
that draws such a bar still counts \file{yes} for X2. Deviation 32 of the survey, written before any reading,
adds an exploratory criterion, X8: is the reported uncertainty computed over random seeds, that is,
over independent training runs? Six fresh AI agents, who had read no paper of the survey, shared the reading
of the 74 papers that are \file{yes} for X2. Two more reread 20 of them blind. Each of the six read papers from
both years, which removes the confound
between readers and year for this criterion alone. Among papers with experiments, X8 goes from 35\% (17 of 48)
to 42\% (21 of 50): $+7$ points, $[-12,\,+25]$, $p = 0.54$ (\cref{tab:x8,fig:x2-x8}). The rise of X2 comes with
no detectable rise of uncertainty computed over random seeds. Error bars spread while their source changed in
no visible way. In 2026, 14 of the 44 papers that draw a bar do not say what it is computed over. X8 was
decided after the main results and read by agents other than the primary readers: it qualifies the X2 verdict
without replacing it.

\fig{x2-x8}{One square per paper with experiments, 48 in 2018 and 50 in 2026, coloured by what its
uncertainty is computed over: random seeds, seeds and episodes pooled, episodes, test items, another source,
an unstated source, no uncertainty. Papers that report an uncertainty go from 63\% to 88\% (X2, a difference that
passes the Holm correction); those that compute it over random seeds go only from 35\% to 42\% (X8, $+7$
points, $[-12,\,+25]$, no detectable change). In 2026, 14 of the 44 papers that draw an error bar do not say
what it is computed over. Data: \file{study/x8\_stats.json}, \file{study/extension\_stats.json}
(\cref{tab:x8}).}

\gentable{x8}

\paragraph{Two intervals that avoid zero without passing the correction}
How the hyperparameters were chosen (Q6) goes from 26\% (12 of 47) to 48\% (24 of 50): $+22$ points,
$[+3,\,+39]$, $p = 0.03$ before correction, $p_{\text{Holm}} = 0.24$. Statistical tools later than 2018
(Y12) go from 2\% (1 of 48) to 14\% (7 of 50): $+12$ points, $[+1,\,+24]$, $p_{\text{Holm}} = 0.66$.
Taken alone, these two intervals avoid zero. With 11 and 12 criteria tested at once, they do not pass
the adjusted threshold. We therefore detect no change after correction, without being able to rule
one out.

\paragraph{No detectable change for the other criteria}
The 16 other criteria have intervals that contain zero. Those that measure a flaw fall into three groups.
\begin{itemize}
  \item \textbf{Rises the survey can neither establish nor rule out.} The method's hyperparameters (Q3,
  55\% then 72\%), held-out evaluation (Q7, 33\% then 45\%) and the generalisation test (Y8, 23\% then
  34\%) climb by more than ten points, as does the choice of hyperparameters (Q6).
  \item \textbf{Flat criteria.} The baselines' hyperparameters (Q4, 35\% then 29\%), baselines tuned with
  the same care (X4, 13\% then 16\%), the measure defined (Y1, 79\% then 78\%), the best runs kept (X3,
  19\% then 14\%), the origin of the baselines' code (Y5, 46\% then 41\%) and the benchmark versions (Y7,
  23\% then 26\%) move by seven points at most. The tuning of baselines, in particular, shows no sign of
  progress.
  \item \textbf{Rare practices.} The statistical test (Q8, 6\% then 12\%) and the library versions (Y10,
  2\% then 8\%) start so low that a rise of six points stays within the noise of the draw.
\end{itemize}

\paragraph{The power of the survey}
Each Newcombe interval spans twenty to forty points. Deviation 30 (b), written before the repository's
computation, puts a number on what the survey could see (\cref{tab:difference-detectable}). Starting from the
2018 share and the actual sample sizes, the two-sided Fisher exact test detects with an 80\% chance only a rise
of 18 to 30 points at the 5\% level, depending on the criterion. At the Bonferroni level over eleven criteria,
it needs 26 to 39 points. Q7 needs 29 points, 37 after correction. The rise actually seen, 12 points, had
only an 18\% chance of being detected at the 5\% level, 4\% after correction. The four established changes
all exceed their detectable difference. X2 exceeds it narrowly: the power of the test for the difference seen falls to
52\% at the Bonferroni level. \enquote{No detectable change} therefore does not mean \enquote{no change}. For
Q3, for instance, the observed difference is $+17$ points, with an interval from $-2$ to $+34$: the data are
compatible with a rise of thirty points as well as with stagnation.

\gentable{difference-detectable}

\section{Rereading flaw by flaw}\label{sec:sr-failles}

This section takes up the results in the order of the flaws of Table~\ref{tab:failles}. The 2018 and
2026 figures are those of the main comparison, unless the first draw is mentioned.

\paragraph{F1, the code}
Linking the code goes from 31\% to 74\% of papers. The flaw is heeded. The official report of the
NeurIPS 2019 reproducibility programme points the same way: the share of accepted papers with code rose
from under 50\% to 74.4\% \parencite{pineau2021improving}. Across the whole first draw, large language
models included, the share is 67\% (30 of 45). One classical RL paper in four still publishes its
result without its code. The rise holds when both years are read on the same version of the paper
(Section~\ref{sec:sr-robustesse}).

\paragraph{F2 and F3, the plan for rerunning and the hardware}
The compute is quantified in 58\% of 2026 papers against 27\% in 2018. All of the method's
hyperparameters are given in 72\% of papers against 55\%, with no detectable change
($p_{\text{Holm}} = 0.58$). Library versions are still rare: 4 classical RL papers out of 50 give them
in 2026, against 1 out of 48 in 2018 (Y10). Across the whole first draw, one paper out of 45 gives
them. Both flaws are partly heeded: compute is quantified, software versions are left unstated.

\paragraph{F4 and F5, the method's hyperparameters and how they were chosen}
The level of Q3 shows no detectable change. How the hyperparameters were chosen is described in 48\%
of 2026 papers against 26\% in 2018, a difference that does not pass the correction. Even in 2026, fewer than
one paper in two says how it found its settings.

\paragraph{F6, the sensitivity to settings}
At least one hyperparameter of the method is swept, with the result shown, in 80\% of 2026 papers
against 44\% in 2018. This flaw was heeded too. The first draw, large models included, gives a similar
level: 32 papers out of 45.

\paragraph{F7, the baselines}
All the baselines' hyperparameters are given in 29\% of 2026 papers (14 of 49), against 35\% in 2018
(17 of 48). The papers that say they tuned their baselines with the same care as their method go from
13\% to 16\%. The practice is rare in both years, with no detectable change. Against Pineau's
60\%, the decline looked sharp (Section~\ref{sec:sr-pineau}). Two facts explain part of it: our grid,
which finds 35\% in 2018, and the papers on language models, which rarely describe their baselines, so
that the whole first draw falls to 18\% (8 of 44).

\paragraph{F8, the performance measure}
The measure is defined in 78\% of 2026 papers against 79\% in 2018 (Y1). When error bars appear, their
meaning is defined in 38 papers out of 44 in 2026 against 25 out of 30 in 2018 (Y2). The measure is
usually defined, in both years, with no detectable change. Y1 is our most fragile criterion: on the
first draw, two readers agreed on only 60\% of papers before its rules were clarified
(Section~\ref{sec:sr-accord}).

\paragraph{F9, the number of runs}
Of the 41 experimental RL papers of the first draw (\file{llm} and \file{classic} families), 20 do not
say how many independent training runs support their results (X1). Four have only one. The median of
the others is five, the most frequent value in the table Pineau showed in 2018. This flaw was not heeded.
Chapter~\ref{chap:06-experiences-resultats} measures with experiment E5 what five random seeds can
settle.

\paragraph{F10, the best runs}
The selection of the best runs shows in 14\% of 2026 papers against 19\% in 2018 (X3), with no
detectable change. In the first draw, 2 papers out of 45 let it show. A paper that keeps its best runs
without saying or showing it escapes our reading: these shares are lower bounds.

\paragraph{F12, the uncertainty and the test}
The flaw splits in two. Reporting an uncertainty has become common: 88\% of 2026 papers against 63\% in
2018. The type of uncertainty is named in 41 of the 44 papers from 2026 that report one, against 25 of
30 in 2018 (Y11). Naming the type of uncertainty does not yet say what it is computed over. Among papers
that draw a bar, the share that compute it over random seeds stays close to one in two, 57\% in 2018 and
48\% in 2026 (X8, \cref{tab:x8}): the lesson of flaw F11 and of experiment B1 has yet to be heard. The
statistical test shows no detectable change: 12\% of 2026 papers against 6\% in
2018 (Q8). The tools that appeared after 2018, such as the IQM or the rliable library, go from 2\% to
14\% (Y12), without passing the correction. Papers draw error bars, but they rarely test.

\paragraph{F13 and F14, the tasks and the implementations}
The main conclusion rests on several tasks in 41 papers out of 50 in 2026, against 43 out of 48 in 2018
(Y3). Among the papers with several tasks, 38 out of 41 show at least one task where their method does
not win, against 34 out of 43 in 2018 (Y4). The origin of the baselines' code is stated in 41\% of 2026
papers against 46\% in 2018 (Y5). The three criteria show no detectable change. The first two
practices were already common in 2018, the third is still a minority practice. In the first draw, only
25 papers out of 45 have several tasks: many papers on language models train on one dataset.

\paragraph{F15, held-out evaluation}
With the same grid on both sides, classical RL goes from 33\% to 45\% for Q7 and from 23\% to 34\% for
the stricter generalisation test (Y8). Both intervals contain zero. The strong apparent rise against
Pineau's 10\% comes from the grid and from the arrival of the papers on language models (Sections~\ref{sec:sr-pineau}
and~\ref{sec:sr-premier}).

\paragraph{F16, the benchmarks}
The version of each benchmark is given in 26\% of 2026 papers against 23\% in 2018 (Y7), with no
detectable change. In the first draw, 16 papers out of the 27 of the \file{llm} family try only one
family of language models (X6).

\section{What withstands the robustness analyses}\label{sec:sr-robustesse}

Table~\ref{tab:robustesse} gathers each alternative analysis and states, for each verdict, whether the
verdict survives it.

\gentable{robustesse}

\paragraph{The provenance of the 2026 papers (a1)}
Within 2026, the 14 papers of the first draw and the 36 of the A3 queue differ on none of the four
criteria (Fisher, $p = 0.73$ for Q5, 1 for Y9 and X7, $0.17$ for X2). Redone with each provenance
alone against 2018, the comparison keeps its sign and an interval that excludes zero in all eight
cases. For Q5: $+47$ points $[+18,\,+65]$ with the first draw, $+41$ $[+20,\,+57]$ with the queue. For
X2, the interval of the queue alone comes close to zero: $+21$ points $[+1,\,+38]$.

\paragraph{Composition by conference (a3)}
The Cochran-Mantel-Haenszel test, stratified by conference, keeps all four verdicts:
$p = 5 \times 10^{-5}$ for Q5, $0.001$ for Y9, $0.005$ for X7, $0.006$ for X2. It supports neither Q6
($p = 0.06$) nor Y12 ($p = 0.15$).

\paragraph{The PDF source (a2)}
This cut sets 45 papers from 2026 read on arXiv against 5 read on a conference site: it settles
nothing, and we do not rely on it. The question it aimed at calls for a direct measurement. An arXiv
version is sometimes longer or more recent than the proceedings version, while the 2018 papers were read in
the proceedings. For Q5, that measurement has been made.

\paragraph{Q5 reread on the same version in both years (deviation 31)}
Code linked is the criterion most exposed to the asymmetry of versions: an author sometimes adds the address
of the repository to an arXiv version posted after the conference. Deviation 31, written before any download,
rereads Q5 in three ways that treat both years alike:
\begin{itemize}
  \item \textbf{both years on the proceedings}: the 45 papers of 2026 read on arXiv were all found in the
  proceedings;
  \item \textbf{both years on the latest arXiv version}: 27 of the 33 papers of 2018 that were \file{no}
  were found there, the other six keep their reading;
  \item \textbf{code available at the conference}: a \file{yes} whose GitHub repositories were all created
  after the last day of the conference becomes \file{no}.
\end{itemize}
A fresh arbiter settled the doubtful cases blind, with years hidden and identifiers replaced. Each reading also
gets a bound where uncertain cases count \file{yes}. The bound of the last one also counts undated
repositories as \file{no}. The verdict holds in all seven readings, the published one, the three rereadings and their three
bounds (\cref{tab:q5-versions,fig:q5-lectures}):
$+41$ points $[+21,\,+56]$ on the proceedings, $+34$ $[+15,\,+51]$ on the latest arXiv version, $+37$
$[+17,\,+52]$ for code available at the conference. The least favourable reading, the latest arXiv version with
uncertain cases counted \file{yes}, still gives $+30$ points $[+11,\,+47]$, $p_{\text{Holm}} = 0.035$. The
asymmetry exists: in 2018, four papers link their code only in their latest arXiv version (C19, C216, C22,
C39). It affects four to six papers out of 48 and explains only a small part of the rise. One repository
alone was created after its conference, that of C219, a year after NeurIPS 2018. One caveat: the arbiter
counts C19 \file{yes} in the arXiv reading, whereas the audit of deviation 19 had visited its page and judged that it led
only to a library released after the paper. Without C19, that reading gives $+37$ points.

\gentable{q5-versions}

\fig{q5-lectures}{Q5, code linked, reread on the same kind of version in both years: the proceedings, the
latest arXiv version, the code available at the conference, each with its bound. Left, the shares of 2018 and
2026, with Pineau's 55\% dotted. Right, the difference and its 95\% Newcombe interval. The difference stays between $+30$ and
$+43$ points and passes the Holm correction in every reading. Our reading of 2018 stays below Pineau's 55\% in
all of them. Data: \file{study/q5\_versions.json}, \file{study/results\_final.json} (\cref{tab:q5-versions}).}

\paragraph{One reader removed at a time (a4)}
Removing each of the sixteen primary readers in turn makes none of the four verdicts fragile. The
weakest lower bound of the Newcombe interval is $+16$ points for Q5 (without H1 or without H3), $+10$
for Y9 (without H2), $+7$ for X7 (without E, K2 or K5) and $+4$ for X2 (without H4). Q6 becomes fragile
without reader H2 ($+18$ points, $[-3,\,+36]$). Y12 is fragile in 8 of the 16 cuts.

\paragraph{Without the thirteen papers read partly by another model}
Removing the thirteen papers finished by \file{claude-opus-5} (deviation 29) leaves 48 classical RL
papers in 2026 and the 50 papers of 2018. Code (Q5, $p_{\text{Holm}} = 0.0001$),
sensitivity (Y9, $0.008$) and compute (X7, $0.04$) still pass the correction. Reported uncertainty no
longer passes it: $+25$ points, $[+8,\,+41]$, $p_{\text{Holm}} = 0.07$. Neither does Q2
($p_{\text{Holm}} = 0.07$). The trends of the trajectory do not change (Section~\ref{sec:sr-trajectoire}).

\paragraph{The search against the \enquote{no} decisions}
On the seven criteria checked, the two samples hold 377 final \file{no} decisions. For 94 of them, the
patterns fixed in advance find no match in the text. The other 283 went to the arbiter. The arbiter
turns 10 of them: 2 in 2018 (C19 for Q5, C120 for Y9) and 8 in 2026 (K058, K063 and LR03 for Q4, K002
for Q7, ML06 for Q8, K056, LR03 and LR09 for X7). The arbiter judges 24 of them doubtful. Redone with the
turned decisions, the main comparison keeps all its verdicts (\file{study/falsify\_no.json}). The difference
in quantified compute grows, to $+37$ points $[+17,\,+53]$. The difference in the baselines' hyperparameters
falls to $-1$ point. If the 24 doubtful decisions are turned as well, reported uncertainty and
sensitivity move just above the threshold ($p_{\text{Holm}} = 0.07$ and $0.054$). Deviation 27 did not
announce this second bound.

\paragraph{Robustness summary}
Code and compute pass the correction in all our checks. Code also passes it when both years are read on the
same version. Sensitivity to settings passes it in all but the
bound of the doubtful decisions, where it stops at $p_{\text{Holm}} = 0.054$. Reported uncertainty fails
in two checks, without the thirteen papers and with the bound of the doubtful decisions, with
$p_{\text{Holm}} = 0.07$ each time. We keep it among the four changes, with this reservation.

\paragraph{The cost and the risk of each practice}
The three raters of the cost and risk of each practice (Section~\ref{sec:sm-robustesse}) agree with
one another: their scores correlate from 0.78 to 0.91 for cost, from 0.90
to 0.93 for risk (Spearman, per pair). Their scores do not predict the observed differences. The Spearman
correlation is $+0.16$ for cost ($p = 0.52$ by permutation) and $-0.04$ for risk ($p = 0.89$). Sweeping
a hyperparameter (Y9), scored costly (median 4), progressed. The versions of benchmarks (Y7) and of
dependencies (Y10), scored nearly free (median 1), did not move. The reading \enquote{heeded where it
cost little} is withdrawn.

\section{Our reading of 2018 against Pineau's figures}\label{sec:sr-pineau}

\paragraph{The difference between the two readings}
The secondary comparison of deviation 12 sets our reading of 2018 against Pineau's figures, on the six
criteria Q3 to Q8 for which she gives a comparable one, with a Holm correction over these six criteria
(\file{extension\_stats.json}, field \file{ours2018\_vs\_pineau}).
\begin{itemize}
  \item \textbf{Q3, the method's hyperparameters}: 55\% (26 of 47) in our reading, 90\% in Pineau's.
  Difference of $-35$ points, $[-50,\,-17]$, $p_{\text{Holm}} = 0.001$.
  \item \textbf{Q7, held-out evaluation}: 33\% (16 of 48) in our reading, 10\% in Pineau's. Difference of
  $+23$ points, $[+7,\,+39]$, $p_{\text{Holm}} = 0.03$.
  \item \textbf{Q4, the baselines' hyperparameters}: 35\% against 60\%. Difference of $-25$ points,
  $[-42,\,-5]$, $p_{\text{Holm}} = 0.06$.
  \item \textbf{Q5, code linked}: 31\% against 55\%. Difference of $-25$ points, $[-42,\,-5]$,
  $p_{\text{Holm}} = 0.06$.
  \item \textbf{Q6 and Q8}: 26\% against 20\% and 6\% against 5\%, with no detectable difference.
\end{itemize}

\gentable{sondage-pineau-contexte}

The first 12 papers of the calibration already gave the same sign: Q3 at 6 of 10, Q7 at 4 of 11
(\file{results\_final.json}). Our grid is stricter than Pineau's on what the paper gives, more lenient
on what counts as held out.

\paragraph{Why our reading differs}
Pineau's draw and her way of settling doubtful cases remain unknown: the slide announces 50 RL papers
from 2018 published at NeurIPS, ICML and ICLR, without saying how they were chosen or how each
criterion was decided. We therefore cannot attribute the difference to one cause. Three mechanisms make
it plausible, none of them demonstrated.
\begin{itemize}
  \item \textbf{Stricter written rules on what is given.} Our rules for Q3 and Q4 require the
  hyperparameters of all the main methods. \enquote{Default values} without a list count only if the
  library version is given. For Q5, \enquote{will be released} and a project page that does not lead
  to the code count as \file{no}. A reading that accepts a partial table or a promise of release gives
  higher shares.
  \item \textbf{A more lenient rule on what is held out.} Our Q7 counts environment random seeds unseen
  in training. Pineau probably had a sharper separation in mind, since her NeurIPS roadmap asks for
  separate tasks and for richer simulators.
  \item \textbf{Another draw.} Our 2018 sample is drawn at random from the 131 papers whose title
  matches our expression. Pineau's sample may contain papers our filter does not see, or exclude
  some.
\end{itemize}
A fourth mechanism was measured for Q5: the version read. Reread on the latest arXiv version, our reading of
2018 goes from 31\% to 40\% (19 of 48), 44\% when uncertain cases are counted (21 of 48,
\cref{tab:q5-versions}). The version explains part of the difference from Pineau's 55\%, not all of it. After
correction, the differences for Q4 and Q5 remain compatible with chance ($p_{\text{Holm}} = 0.06$). Those for Q3 and
Q7 do not.

\paragraph{The verdict rule that follows}
Comparing our 2026 papers with Pineau's figures would mix the change in practice with the change of
grid. Deviation 20 draws the verdict rule from this (Section~\ref{sec:sm-verdict}). Pineau's figures
stay on display as context, in Section~\ref{sec:sr-premier}.

\paragraph{What the original rule would have given}
Deviation 30 (c), written before the repository's computation but after a sketch by the technical review
(Section~\ref{sec:ecarts-audit}), applies to the same data the original rule of deviation 7, which
deviation 20 replaced: each of Pineau's criteria was judged against her figures, then against our
reading of 2018 (Section~\ref{sec:sm-verdict}, \cref{tab:regle-origine}). In classical RL in 2026, no criterion would have been \enquote{changed and
confirmed}. Q4, Q6 and Q7 would have \enquote{changed} against Pineau, without confirmation against our 2018.
Q5, our strongest verdict, would have been \enquote{ambiguous}: it passes the correction against our reading of
2018, not against Pineau's figures ($+18$ points, $p_{\text{Holm}} = 0.19$). In the first draw, large models
included, only Q7 would have been \enquote{changed and confirmed}, driven by the papers on language models
(Section~\ref{sec:sr-premier}). The original rule says nothing about X2, X7 and Y9, for which Pineau gives no
figure. This table does not overturn the decision of deviation 20, which rests on the calibration difference measured
for Q3, 55\% in our reading against 90\% in Pineau's. It makes that decision checkable.

\gentable{regle-origine}

\section{The first draw and the language models}\label{sec:sr-premier}

\paragraph{The first draw against Pineau}
Figure~\ref{fig:remesure-dumbbell} compares the 45 papers with experiments of the first draw, large
language models included, with Pineau's figures. Figure~\ref{fig:ecarts-pineau} draws the differences with
their Newcombe intervals, Figure~\ref{fig:remesure-ic} the 2026 shares with their Wilson intervals.
Three intervals stand clearly away from zero. The method's hyperparameters decline, from 90\% to 67\%
($-23$ points, $[-39,\,-7]$, $p_{\text{Holm}} = 0.02$). Those of the baselines decline, from 60\% to
18\% ($-42$ points, $[-57,\,-22]$, $p_{\text{Holm}} = 0.0002$). Held-out evaluation rises, from 10\% to
70\% ($+60$ points, $[+42,\,+73]$). Code goes from 55\% to 67\% with no detectable difference. The statistical
test goes from 5\% to 4\%. Two causes add up in them, the change in practice and the difference between the grids.
These differences serve as context, never as a verdict.

\fig{remesure-dumbbell}{Pineau's grid, applied in 2018 and in 2026, as context. Each row is a reading
criterion. The 2018 dot gives the share of the 50 papers read by Pineau that meet the criterion. The
2026 dot gives the share among the 45 papers with experiments of our first draw from ICLR 2026, ICML
2026 and NeurIPS 2025, large language models included. The two dots do not come from the same grid:
our reading of 2018 finds 55\% for the method's hyperparameters, where Pineau found 90\%. The
statistical test, left out of the drawing, goes from 5\% to 4\%. Data:
\file{study/audit\_stats.json}.}

\fig{ecarts-pineau}{Difference between the share of the 45 papers with experiments of our first 2026 draw and
that of the 50 papers from 2018 read by Pineau, criterion by criterion, in percentage points. The dot
is the observed difference, the bar its 95\% Newcombe interval, which accounts for the uncertainty of both
samples. Above zero, the practice is more frequent in 2026. The difference mixes the change in practice with
the difference between our grid and Pineau's: it serves as context. The verdicts come from the
same-grid comparison. Data: \file{study/audit\_stats.json}.}

\fig{remesure-ic}{Share of the 45 papers with experiments of the first 2026 draw that meet each
criterion, with its 95\% Wilson interval. The draw mixes classical RL with papers on large language
models. With a sample of this size, a difference of less than fifteen points from 2018 cannot be told apart
from the luck of the draw. Data: \file{study/results\_final.json}.}

\paragraph{The papers on language models}
Table~\ref{tab:sondage-lrl} sets, within the first draw, the 27 papers of the \file{llm} family against
the 14 classical RL papers, which several criteria clearly separate. Held-out evaluation is almost universal in the
\file{llm} family (26 of 27) and a minority practice in the classical RL of the first draw (5 of 13):
collections of mathematics problems separate training and test problems by construction. Uncertainty is
rare in the \file{llm} family (4 of 27) and universal in the classical RL of the first draw (14 of 14).
The baselines' hyperparameters are all given in 5 \file{llm} papers out of 27. Sixteen \file{llm}
papers out of 27 try only one family of base models. The apparent rise of held-out evaluation against
Pineau comes almost entirely from these papers. So does the apparent decline of uncertainty in the
first draw.

\gentable{sondage-lrl}

\paragraph{A change of population}
The \file{llm} family did not exist in 2018. A difference between the first 2026 draw and 2018 that runs
through these papers first measures the arrival of a new population, with its own habits. It says
nothing about the change in practice within a population that stayed comparable. The verdicts of this
report therefore hold for classical RL alone.

\section{The trajectory from 2018 to 2026}\label{sec:sr-trajectoire}

\paragraph{The shares year by year}
Table~\ref{tab:trajectoire} and Figure~\ref{fig:trajectoire} give, for each practice, the share of
classical RL papers that follow it in 2018 (48 papers), in 2019, 2020, 2022 and 2024 (20 papers per
year, ICML and NeurIPS) and in 2026 (50 papers). Code linked goes from 31\% in 2018 to 60\%, 50\%, 70\%
and 75\% from 2019 to 2024, then 74\% in 2026. Sensitivity to a hyperparameter is at 44\%, then 35\%,
45\%, 60\%, 85\% and 80\%. Compute quantified is at 27\%, then 30\%, 40\%, 40\%, 55\% and 58\%. Reported
uncertainty is at 63\%, then 80\%, 85\%, 95\%, 100\% and 88\%. The statistical test is at 6\%, then
5\%, 10\%, 10\%, 25\% and 12\%. With 20 papers, a share is uncertain by about twenty points either way.

\gentable{trajectoire}

\fig{trajectoire}{Share of classical RL papers that follow each practice, year by year: 48 papers from
2018 (those with experiments), 20 per year for 2019, 2020, 2022 and 2024 drawn from ICML and NeurIPS,
50 from 2026. The dotted vertical line marks the NeurIPS checklist (2019). Code linked, sensitivity to
a hyperparameter, compute quantified and reported uncertainty climb steadily. The statistical test, in
grey, shows no trend. With 20 papers, a share is uncertain by about twenty points either way. Data:
\file{study/trajectory.json}.}

\paragraph{The trends}
The Cochran-Armitage test, with the years as scores, gives a clear trend for the four practices after
the Holm correction over the five tested criteria: $z = 4.29$ for code
($p_{\text{Holm}} = 7 \times 10^{-5}$), $z = 4.78$ for sensitivity ($9 \times 10^{-6}$), $z = 3.44$ for
compute ($0.002$), $z = 3.39$ for uncertainty ($0.002$). The statistical test shows no trend
($z = 1.53$, $p = 0.13$).

\paragraph{No jump right after 2019}
From 2018 to 2019 and 2020 pooled, no difference passes the correction. Code rises by $+24$ points,
$[+3,\,+42]$, with $p = 0.03$ before correction and $p_{\text{Holm}} = 0.16$ after. Uncertainty rises
by $+20$ points, $[+1,\,+37]$, $p_{\text{Holm}} = 0.23$. Sensitivity does not rise ($-4$ points). From
2019 and 2020 pooled to 2022 and 2024 pooled, only sensitivity rises clearly, from 40\% to 73\%, with an
interval from $+11$ to $+50$ points and $p_{\text{Holm}} = 0.03$. These series describe a slow rise that
began before the checklist or with it. They do not say what caused it. \textcite{coakley2026shift}
likewise see the share of released code rise before the checklists, across five artificial
intelligence conferences.

\paragraph{Without the thirteen papers}
With the papers finished by \file{claude-opus-5} removed, the four trends still pass the correction
($p_{\text{Holm}} < 0.004$) and the statistical test still shows none. The later rise of sensitivity
keeps $p_{\text{Holm}} = 0.02$ (\file{trajectory\_sans\_opus5.json}).

\paragraph{The statistical test's peak in 2024}
The statistical test reaches 25\% of papers in 2024 (5 of 20), then falls back to 12\% in 2026 (6 of
50). With 20 papers, this peak stays within the margin of uncertainty. We do not read it as a fact.

\section{Agreement between readers}\label{sec:sr-accord}

Table~\ref{tab:accord} and Figure~\ref{fig:accord-lecteurs} give the agreement of the blind double
readings, criterion by criterion, on the raw readings before harmonisation.

\gentable{accord}

\fig{accord-lecteurs}{Agreement between the blind double readings, criterion by criterion, on the raw
readings, before harmonisation: share of identical decisions and Cohen's kappa when both readers vary.
Four pairs of readings of ten papers each: G (grid v1) and G2 (grid v2) on the first 2026 draw, R2 on
the added papers of 2018 and 2026, R3 on the trajectory from 2019 to 2024. The 70\% agreement
threshold, below which a criterion is called fragile, is marked. Data: \file{study/agree\_v2.json},
\file{study/audit\_stats.json}.}

\paragraph{Thirty papers reread, with their intervals}
The four pairs of double readings cover 30 distinct papers. G and G2 reread the same ten papers from 2026, one
with grid v1, the other with grid v2. R2 rereads ten added papers, five from 2018 and five from 2026. R3 rereads
ten papers of the trajectory. Deviation 30 (a), written before this computation, pools these readings and
gives each share of agreement its Wilson interval, each kappa a percentile bootstrap interval over papers
(\cref{tab:accord-intervalles}). On the four criteria of the verdicts, the 120 doubled decisions are identical:
30 of 30 for each criterion, with a Wilson interval from 89\% to 100\%. Three criteria keep a kappa below 0.70:
the measure defined (Y1), the origin of the baselines' code (Y5) and the exclusions stated (Y6). One pair alone
says little: on ten papers, 9 agreements have the interval [60, 98]\%. G's kappa of 0.80 on Q3 has the
interval [0.31, 1]. Agreement on 2018 papers rests on the five papers of R2: 5 agreements out of 5 have the
Wilson interval [57, 100]\%.

\gentable{accord-intervalles}

\paragraph{Grid v1 on the first draw}
The second reader G and the primary readers agree on 127 decisions out of 130 (98\%), that is, 13
criteria on 10 papers. The three disagreements concern Q3, Q7 and X4, with kappas of 0.80, 0.78 and
0.74. The kappa is 1 for the other criteria on which both readers vary.

\paragraph{Grid v2 on the first draw}
On the same ten papers, G2 and the primary readers agree completely on 9 criteria out of 12. Agreement
falls to 90\% on exclusions (Y6, kappa 0.80), to 80\% on the implementation of the baselines (Y5, kappa
0.52) and to 60\% on the performance measure (Y1, kappa 0.31). Y1 is the only fragile criterion in all
the double readings. Deviations 17 and 19 then clarified its rules.

\paragraph{The added papers}
R2 rereads ten papers drawn from the additions of 2018 and 2026 with both grids. Agreement reaches at
least 80\% on each of the 23 criteria and is complete on 17 of them. The lowest agreements are those of
Y4 (80\%, kappa 0.52) and Y6 (80\%, kappa 0.41).

\paragraph{The trajectory}
R3 rereads ten papers of A4, with the same outcome: none of the 23 criteria falls below 80\% agreement
and 17 stay complete, the five tested criteria of the trajectory among them (Q5, Y9, X7, X2, Q8). The lowest
kappas are those of Q7 and Y8 (80\%, kappa 0.41).

\paragraph{The stability of one model}
All the readers are copies of the same model, except the two that finished thirteen papers on
\file{claude-opus-5} (Section~\ref{sec:sm-lecteurs}). The agreement of the double readings therefore
measures how stably this one model reads (Section~\ref{sec:sm-doubles}). A human might have read the same
passages differently. The quotations and pages in \file{study/codes/} and \file{study/codes\_v2/} let
anyone check each decision. A reader from another family of models, \file{gemma3:12b}, rereads the 98 papers
of both samples on the criteria of the verdicts (deviation 33, written before this reading). Its reading was
not finished when this report was compiled and does not appear in it.

\paragraph{Year and readers disentangled}\label{sec:sr-annees}
In the main comparison, the readers of 2018 are not those of 2026 (Section~\ref{sec:sm-annees}).
Deviation 36 has the 98 experimental papers of that comparison reread by eight fresh agents, T1 to T8,
each of whom reads six papers from 2018 and six or seven from 2026 (Section~\ref{sec:sm-doubles}). With
these readings alone, the share of \enquote{yes} goes from 31\% to 72\% for code linked (Q5), from 42\%
to 78\% for sensitivity to a hyperparameter (Y9), from 31\% to 66\% for quantified compute (X7) and from
69\% to 86\% for reported uncertainty (X2) (\cref{tab:r36}).

\gentable{r36}

The logistic mixed model gives each reader its own leaning towards \enquote{yes}, then measures the
year effect for a given reader. The standard deviation $s$ of the reader effect is 0 for Q5 and Y9, 0.24 for X7 and 0.44
for X2: none of them stands out from zero ($p \geq 0.25$). The leaning towards \enquote{yes} therefore
varies little from one reader to the next. After the Holm correction over the four criteria, Q5, Y9
and X7 pass both tests fixed in advance, Fisher's test on the difference and the test of the year
effect in the model. Their verdicts hold. For X2, the difference is $+17$ points, $[+0.6,\,+33.0]$.
The year effect passes the correction ($p_{\text{Holm}} = 0.037$), but Fisher's test narrowly misses it
($p_{\text{Holm}} = 0.053$), so the rule calls this verdict \enquote{same direction, not established}.
X2 was already the most fragile verdict (Section~\ref{sec:sr-robustesse}). Every Newcombe interval
overlaps the one of the main comparison. In the reading least favourable to Q5, where every project
page counts as \enquote{yes}, its rise still passes the Holm correction, from 42\% to 74\%
($p_{\text{Holm}} = 0.004$).

\paragraph{The new readings match the original codes}
Each paper thus has two complete readings: a fresh reader's and the original final code, harmonised.
They agree on 94\% to 99\% of the papers depending on the criterion, on 90\% to 100\% depending on the
criterion and the year (\cref{tab:r36-accord}). No criterion shows a clearly different agreement in
2018 and 2026: the smallest Fisher $p$ is 0.11. For X7, all six disagreements go towards
\enquote{yes}, two in 2018 and four in 2026. Eight X7 \enquote{yes} rest on measured durations
(running, inference or training times) rather than on a stated total compute. Deviation 19 leaves
this borderline case to the reader's judgement. Three of these papers are from 2018 and five from
2026. Six of the eight are already \enquote{yes} in the original code.
Chapter~\ref{chap:10-limites} states what this rereading does not separate.

\gentable{r36-accord}

\section{The declarations of the papers}\label{sec:sr-declarations}

\paragraph{Statistical significance}
The NeurIPS checklist question on statistical significance settles for error bars \enquote{or other
appropriate information}. Of the 30 NeurIPS 2025 papers read, 18 answer \enquote{Yes} to it. None of
the 18 applies a test: the checklist does not ask for the test and the papers do not run one. One of
the 18 reports no uncertainty either. The others report an uncertainty and stop there
(\file{study/checklist.json}).

\paragraph{Code, settings and compute}
On code, 22 papers answer \enquote{Yes}, 5 of them without a link to their code in the version we
read, or 23\% $[10,\,43]$. Four have no link even in the proceedings version. Three write in their own
justification that the code will be released later. On settings, 22 papers answer
\enquote{Yes}, 6 of them without giving all the hyperparameters of their method according to our rule.
On compute, 21 papers answer \enquote{Yes}, 6 of them without quantifying any compute time, which the
question explicitly asks for. One paper has no readable checklist in the proceedings. A ticked box does
not guarantee the practice it names.

\section{Sensitivity to the title filter}\label{sec:sr-filtre}

\paragraph{The RL papers the filter misses}
Of the 400 papers drawn outside the title filter, 25 have an abstract that matches the regular
expression: 6.2\%, with a Wilson interval from 4.3\% to 9.1\% (\file{study/title\_filter.json}).
Scaled to the population outside the filter, this share gives about 1,037 missed RL papers, from
708 to 1,503, or 45\% of all RL papers according to this criterion (from 36\% to 55\%). This figure
is an upper bound: the rule is mechanical and also counts false positives, such as a paper where
\enquote{RL} stands for representation learning. Most of the papers found deal with large language
models.

\paragraph{The twenty hidden-RL papers}
Read with both grids, the first 20 hidden-RL papers resemble the first 2026 draw: code linked in 60\%
of them (67\% in the draw), an uncertainty in 50\% (47\%). Thirteen deal with large language models.
The 7 classical RL papers are too few to test our verdicts. The filter misses many papers, above all on
language models, with no sign that these papers follow different practices.

\section{Synthesis}\label{sec:sr-synthese}

With the same grid on both sides, in classical RL, four practices changed between 2018 and 2026:
linking the code, showing the sensitivity to a hyperparameter, quantifying the compute, reporting an
uncertainty. Code and compute withstand every cut, sensitivity nearly every one. Reported uncertainty
fails in two checks, with $p_{\text{Holm}} = 0.07$. Eight fresh readers, each reading both years, find
the rises in code, sensitivity and compute again. The rise in uncertainty keeps its direction without
being established. Code linked also keeps its verdict when both years are read on the same version.
All four practices rise slowly from 2018 to 2024, with no jump right after the 2019 checklist.

The statistical test, baselines described and tuned with care, held-out evaluation, benchmark versions
and library versions show no detectable change. The tuning of baselines shows no sign of progress. For
the rest, the survey cannot decide, since it could only see rises of about thirty points. Error bars
spread while the share computed over random seeds moved in no detectable way (X8). The clear changes
concern what accompanies the result: the code link, the declared compute, the effect of a setting, the
error bar.

Our reading of 2018 differs from Pineau's on the method's hyperparameters and on held-out evaluation,
hence the same-grid verdict rule. The original verdict rule, applied to the same data, would have
confirmed no change in classical RL. Chapter~\ref{chap:10-limites} gathers the limitations of these
results. Chapter~\ref{chap:11-conclusion} sets them against the experiments.
